% ============================================================================
% SocialGuard AI — IEEE Format Research Paper
% Compile with: pdflatex SocialGuard_AI_IEEE.tex
% Or upload directly to Overleaf.com
% ============================================================================
\documentclass[conference]{IEEEtran}

% --- Packages ---
\usepackage{cite}
\usepackage{amsmath,amssymb,amsfonts}
\usepackage{algorithmic}
\usepackage{graphicx}
\usepackage{textcomp}
\usepackage{xcolor}
\usepackage{booktabs}
\usepackage{multirow}
\usepackage{hyperref}
\usepackage{tabularx}
\usepackage{float}

\hypersetup{
    colorlinks=true,
    linkcolor=blue,
    citecolor=blue,
    urlcolor=blue
}

\def\BibTeX{{\rm B\kern-.05em{\sc i\kern-.025em b}\kern-.08em
    T\kern-.1667em\lower.7ex\hbox{E}\kern-.125emX}}

\begin{document}

% ============================================================================
% TITLE
% ============================================================================
\title{SocialGuard AI: An Ensemble-Driven, Explainable Detection Pipeline for Identifying Inauthentic Profiles Spanning Six Social Platforms}

% --- AUTHORS (UPDATE WITH YOUR DETAILS) ---
\author{
\IEEEauthorblockN{First Author Name}
\IEEEauthorblockA{
\textit{Department of Computer Science/Engineering}\\
\textit{Your University Name}\\
City, Country\\
email@university.edu}
\and
\IEEEauthorblockN{Second Author Name (Guide/Co-Author)}
\IEEEauthorblockA{
\textit{Department of Computer Science/Engineering}\\
\textit{Your University Name}\\
City, Country\\
email@university.edu}
}

\maketitle

% ============================================================================
% ABSTRACT
% ============================================================================
\begin{abstract}
Inauthentic accounts and automated agents pervade contemporary social networks, undermining collective trust and enabling coordinated manipulation campaigns. We introduce \textbf{SocialGuard AI}, a unified detection framework covering Instagram, Twitter (X), Reddit, Facebook, LinkedIn, and YouTube through platform-tailored behavioral feature extraction and a three-algorithm weighted ensemble --- XGBoost, Random Forest, and Logistic Regression --- whose outputs are fused via AUC-calibrated soft voting. Every prediction is accompanied by SHAP-derived explanations that expose which profile attributes most influence the verdict. The system ships as a containerized full-stack application: a FastAPI backend with lifecycle-managed startup, X-Request-ID tracing, and rate-limit quota headers; a React 19 interface offering dual-theme rendering, side-by-side comparison, batch CSV ingestion, and PDF reporting; a Manifest~V3 Chrome extension; and a key-authenticated public REST endpoint for programmatic consumers. An intelligent chatbot backed by Groq's Llama~3.3-70B (with Gemini failover) guides users through result interpretation. On genuine crowd-sourced datasets the pipeline reaches 98.6\% accuracy for Instagram, 99.29\% for Twitter, 94.0\% for Reddit, and 95.0\% for Facebook, each surpassing 0.97 AUC. System dependability is validated through 110 automated checks spanning backend unit tests, React component tests, and browser-level integration flows, orchestrated by a four-stage continuous-integration pipeline.
\end{abstract}

\begin{IEEEkeywords}
inauthentic account identification, cross-platform analysis, gradient-boosted ensembles, SHAP interpretability, behavioral profiling, browser extension, continuous integration
\end{IEEEkeywords}

% ============================================================================
% I. INTRODUCTION
% ============================================================================
\section{Introduction}

Digital social platforms now connect upward of five billion individuals globally \cite{statista2026}, yet the very openness that fosters communication also invites exploitation. Automated agents, sockpuppet networks, and spam-driven profiles erode the quality of public discourse, facilitate phishing operations, and artificially amplify engagement signals \cite{cresci2020decade, ferrara2016rise}. Independent analyses suggest that roughly one in seven active accounts on leading networks may be inauthentic \cite{varol2017}.

Most prior detection efforts concentrate on a single network, producing siloed classifiers that cannot generalize across the heterogeneous feature landscapes of different platforms. Equally problematic, many systems function as opaque binary oracles --- issuing a ``fake'' or ``genuine'' verdict without exposing the reasoning behind it. Such opacity hampers trust among end-users and limits utility for content moderators who require actionable intelligence.

The work described herein addresses both shortcomings through \textbf{SocialGuard AI}, whose principal contributions are:

\begin{enumerate}
    \item \textbf{Unified six-platform coverage.} A single codebase ingests platform-specific behavioral and structural signals from Instagram, Twitter, Reddit, Facebook, LinkedIn, and YouTube, each through a dedicated feature-engineering pipeline processing between 6 and 16 attributes.
    \item \textbf{Ensemble intelligence with confidence calibration.} Three fundamentally different learners --- gradient-boosted trees, bootstrap-aggregated forests, and a regularized linear model --- are fused through AUC-weighted soft voting. Inter-model agreement serves as an intrinsic confidence gauge.
    \item \textbf{Transparent reasoning via SHAP.} Every prediction is deconstructed into per-feature Shapley contributions, enabling users to inspect precisely which profile characteristics pushed the score toward ``fake'' or ``genuine.''
    \item \textbf{End-to-end deployment readiness.} The system is packaged with Docker Compose containerization, a four-job CI pipeline, dual authentication layers (Firebase tokens and API keys), persistent rate limiting with quota headers, distributed request tracing, and a Chrome browser extension.
    \item \textbf{Comparative and temporal analytics.} A split-screen comparison mode, timeline trend charts, confidence distribution histograms, and weekly activity heatmaps equip analysts with multidimensional situational awareness.
\end{enumerate}

The remainder of this manuscript is structured as follows. Section~II surveys prior art. Section~III delineates the system architecture. Section~IV elaborates on methodology --- datasets, feature construction, and training procedures. Section~V describes implementation specifics. Section~VI presents and interprets experimental outcomes. Section~VII discusses strengths, limitations, and prospective extensions. Section~VIII offers concluding remarks.

% ============================================================================
% II. RELATED WORK
% ============================================================================
\section{Related Work}

\subsection{Detectors Confined to One Network}

Individual-platform classifiers dominate the literature. Varol and colleagues assembled a feature bank exceeding one thousand behavioral dimensions --- spanning temporal regularity, linguistic markers, and network topology --- to produce the BotOrNot scoring service (later relaunched as Botometer) for Twitter \cite{varol2017}. In a follow-up line of inquiry, Yang et al.\ demonstrated that stripping the feature set down to lightweight account metadata still yields competitive separation, opening the door to low-cost deployment on resource-constrained infrastructure \cite{yang2020scalable}. Within the Instagram ecosystem, Akyon and Kalfaoglu trained convolutional architectures on engagement cadence signals to surface counterfeit personas \cite{akyon2019instagram}.

On the Facebook front, Fire et al.\ probed anomalous topologies within friend-request graphs as indicators of deceptive intent \cite{fire2014friend}, while Boshmaf et al.\ mounted a controlled socialbot infiltration campaign demonstrating how trust-based friending policies can be weaponized at scale \cite{boshmaf2011socialbot}. For Reddit, Hurtado et al.\ leveraged temporal posting rhythms coupled with lexical fingerprinting to partition genuine participants from orchestrated agents \cite{hurtado2019bot}.

\subsection{Multi-Network Initiatives}

Cross-boundary detection remains comparatively nascent. Cresci's ten-year retrospective charted the escalating arms race between bot creators and detectors, explicitly identifying unified multi-network architectures as an unaddressed frontier \cite{cresci2020decade}. Kudugunta and Ferrara furnished empirical evidence that behaviorally grounded features can traverse network boundaries when channeled through sufficiently expressive deep models \cite{kudugunta2018deep}. Despite these signposts, no earlier system has shipped a six-platform, ensemble-backed, explainability-equipped production deployment.

\subsection{Explanation-Aware Security Systems}

Shapley-based feature attribution, codified by Lundberg and Lee \cite{lundberg2017shap}, has permeated adversarial-defense pipelines. Rathore et al.\ embedded SHAP within network-perimeter intrusion detectors and measured a statistically significant uplift in operator trust toward automated alerts \cite{rathore2023xai}. Mohawesh et al.\ assessed the broader landscape of content-moderation systems and concluded that transparency has graduated from a desirable property to a deployment prerequisite \cite{mohawesh2023explainability}.

\subsection{Tree-Based Gradient Machines}

Chen and Guestrin's XGBoost \cite{chen2016xgboost} has accumulated a track record of benchmark victories on structured-data tasks spanning bot flagging \cite{yang2020scalable}, financial anomaly surfacing \cite{zhang2019fraud}, and perimeter-defense log analysis \cite{dhaliwal2018effective}. Its tree-native architecture permits exact Shapley computation via the TreeExplainer kernel, rendering it a natural backbone for pipelines that require both high accuracy and interpretable outputs.

% ============================================================================
% III. SYSTEM ARCHITECTURE
% ============================================================================
\section{System Architecture}

SocialGuard AI adopts a four-tier modular layout comprising a client layer, an API gateway, an intelligence core, and a persistence tier, depicted schematically in Fig.~\ref{fig:architecture}.

\begin{figure}[htbp]
\centerline{
\begin{tabular}{|c|}
\hline
\textbf{Presentation Tier} \\
React 19 SPA (Dark/Light Themes) \\
Chrome Extension (MV3) | Comparison View \\
\hline
\textbf{Gateway Tier} \\
FastAPI + Uvicorn, Pydantic v2 Validation \\
X-Request-ID Tracing | X-RateLimit Headers \\
Public API v1 (Key-Gated) \\
\hline
\textbf{Intelligence Tier} \\
Ensemble Pipeline (XGBoost + RF + LR) \\
SHAP TreeExplainer | PDF Engine | AI Chat \\
\hline
\textbf{Persistence Tier} \\
Firebase Auth | Firestore | 17 Serialized Models \\
\hline
\end{tabular}
}
\caption{Four-tier layout of the SocialGuard AI platform showing the gateway, intelligence, and persistence boundaries.}
\label{fig:architecture}
\end{figure}

\subsection{Server-Side Organization}

The Python backend, built atop FastAPI, is partitioned into six cohesive packages:

\begin{itemize}
    \item \textbf{Model Loader (\texttt{ml/engine.py}):} Deserializes 17 serialized classifiers (six XGBoost, five Random Forest, five Logistic Regression, plus ensemble weight metadata), initializes SHAP TreeExplainers, maps feature names per platform, and houses the risk-interpretation logic that downgrades confidence for synthetically trained models.
    \item \textbf{Voting Combiner (\texttt{ml/ensemble.py}):} Computes AUC-weighted soft votes across the three algorithm families, measures inter-model spread, and classifies agreement strength as unanimous, high, moderate, or low.
    \item \textbf{Endpoint Modules (\texttt{routes/}):} Eight self-contained router files cover predictions, batch processing, URL scanning, health diagnostics, model-comparison metrics, PDF/ZIP reporting, chatbot proxying, and the public API surface.
    \item \textbf{Security Stack (\texttt{middleware.py}):} Three stacked middleware layers --- (1) request logging injecting a unique X-Request-ID into every log line and response header, (2) sliding-window rate limiting with file-backed persistence and X-RateLimit-Limit/Remaining/Reset headers, and (3) Firebase ID-token verification with a development-mode bypass.
    \item \textbf{Input Gates (\texttt{schemas.py}):} Pydantic v2 models enforce Field-level constraints (\texttt{ge}, \texttt{le}, \texttt{min\_length}, \texttt{max\_length}) on all 96+ input attributes, rejecting malformed payloads before they reach any model.
    \item \textbf{Lifecycle Manager:} An async context manager hooks into FastAPI's lifespan protocol, performing model-availability audits and emitting a startup banner during initialization, and executing graceful resource cleanup on shutdown.
\end{itemize}

\subsection{Client-Side Organization}

The browser interface, authored in React 19 with Vite 8, comprises roughly twenty components grouped by responsibility:

\begin{itemize}
    \item \textbf{Identity:} Firebase-backed sign-up and sign-in flows supporting email/password, Google OAuth, and GitHub OAuth providers.
    \item \textbf{Detection:} Dynamically rendered forms that adapt input fields to the selected platform configuration, with a toggle enabling ensemble multi-model mode.
    \item \textbf{Comparison:} A split-screen layout allowing simultaneous analysis of two profiles, showing dual gauge visualizations, paired SHAP bar charts, and an auto-generated feature-difference table.
    \item \textbf{Charting:} SVG-based probability gauges, gradient-filled ensemble contribution bars, Recharts-powered SHAP importance displays, timeline area charts, confidence histograms, and weekly activity bar charts.
    \item \textbf{Dashboard:} A four-tab analytics console featuring overview counters, temporal trend lines, model radar/bar comparisons, and threat-density heatmaps.
    \item \textbf{Theming:} A dark/light toggle driven by CSS custom properties, persisted to localStorage, and initialized from the operating system's \texttt{prefers-color-scheme} media query.
    \item \textbf{Developer Console:} Key generation with masked preview, per-key usage statistics, cURL snippet generation, and one-click revocation.
    \item \textbf{Bulk Ingestion:} CSV upload supporting multi-record analysis with aggregate statistics and downloadable ZIP archives of individual PDF reports.
    \item \textbf{Conversational Aid:} A floating chatbot widget backed by Groq's Llama 3.3-70B (primary) and Google Gemini (fallback), proxied entirely server-side to shield API credentials.
\end{itemize}

\subsection{Browser Extension}

A Chrome extension conforming to Manifest V3 enables contextual scanning while a user browses a social profile. It communicates with the backend through the same REST endpoints and renders results in a lightweight popup overlay.

% ============================================================================
% IV. METHODOLOGY
% ============================================================================
\section{Methodology}

\subsection{Training Corpora}

Six labeled collections --- four harvested from public repositories and two purposefully synthesized --- underpin the platform classifiers. Table~\ref{tab:datasets} catalogs their characteristics.

\begin{table}[htbp]
\caption{Training Corpus Characteristics}
\label{tab:datasets}
\centering
\begin{tabular}{lcccc}
\toprule
\textbf{Platform} & \textbf{Origin} & \textbf{Records} & \textbf{Dims.} & \textbf{Balance} \\
\midrule
Instagram & Kaggle\cite{instagram_dataset} & 5{,}000 & 11 & 50:50 \\
Twitter & Kaggle\cite{twitter_dataset} & 2{,}818 & 14 & $\sim$47:53 \\
Reddit & Custom\cite{reddit_dataset} & 500 & 6 & $\sim$44:56 \\
Facebook & Kaggle\cite{facebook_dataset} & 600 & 13 & Skewed \\
LinkedIn & Synthesized & 2{,}000 & 12 & 50:50 \\
YouTube & Synthesized & 2{,}000 & 10 & 50:50 \\
\bottomrule
\end{tabular}
\end{table}

\textbf{Instagram.} Sourced from a Kaggle repository \cite{instagram_dataset}, the collection holds 696 annotated accounts profiled across avatar availability, alphanumeric username composition, biography text presence, publication tally, inbound follower count, and outbound followee count.

\textbf{Twitter.} Amalgamating verified human-operated handles with confirmed bot identities \cite{twitter_dataset}, the 2,818-record corpus spans 14 axes: activity throughput (tweet and favourite tallies), social-graph topology (follower and friend counts, public-list memberships), and configuration metadata (verification seal, stock avatar flag, geolocation toggle, backdrop customization state, timezone displacement, and privacy setting).

\textbf{Reddit.} Originating from a ``dead-internet'' forensic inquiry \cite{reddit_dataset}, this lean 500-record set captures account tenure, aggregate karma, comment-level sentiment polarity, mean token length, and hyperlink prevalence.

\textbf{Facebook.} The Kaggle-hosted corpus \cite{facebook_dataset} encodes group membership count, account chronological age, publication frequency, URL-sharing fraction, media-sharing fraction, and per-post engagement indicators (comment depth, like accumulation, tag density).

\textbf{LinkedIn and YouTube.} Publicly labeled ground-truth datasets do not exist for either network. To validate the framework's extensibility, we fabricated 2{,}000 records per platform from domain-calibrated probability distributions: authentic-profile draws featured richer engagement and prolonged account existence, while fabricated-profile draws exhibited anemic activity and curtailed lifespans. \textbf{The resulting class boundaries are non-overlapping}, producing trivially partitionable data and artificially perfect metrics (100\% across every measure). Outputs from these two classifiers therefore carry a persistent ``Synthetic'' badge within the interface, the REST responses (\texttt{data\_source: "synthetic"}), and the automated test assertions, accompanied by an automatic confidence penalty. They must \textit{not} be treated as evidence of real-world generalization capability.

\subsection{Attribute Construction}

Platform-specific engineering transforms raw inputs into discriminative signals. Table~\ref{tab:features} enumerates the principal derived attributes.

\begin{table}[htbp]
\caption{Key Derived Attributes by Platform}
\label{tab:features}
\centering
\begin{tabular}{ll}
\toprule
\textbf{Platform} & \textbf{Constructed Attributes} \\
\midrule
Instagram & Follower-to-followee ratio, posts-per- \\
          & follower density, completeness index \\
Twitter   & Full 14-dim vector including geo, bg, \\
          & timezone, and protection flags \\
Reddit    & Daily karma accumulation rate \\
Facebook  & Age normalization (years $\to$ days) \\
LinkedIn  & Direct use of all 12 dimensions \\
YouTube   & Direct use of all 10 dimensions \\
\bottomrule
\end{tabular}
\end{table}

For Instagram, three compound signals are computed:

\begin{equation}
R_{ff} = \frac{N_{\text{followers}}}{N_{\text{following}} + 1}
\end{equation}

\begin{equation}
R_{pf} = \frac{N_{\text{posts}}}{N_{\text{followers}} + 1}
\end{equation}

\begin{equation}
S_{\text{profile}} = I_{\text{pic}} + I_{\text{bio}} + I_{\text{name}}
\end{equation}

\noindent where $R_{ff}$ quantifies social reciprocity, $R_{pf}$ measures content-production density, and $S_{\text{profile}}$ aggregates binary completeness indicators.

For Reddit, a karma velocity metric is derived:

\begin{equation}
V_{\text{karma}} = \frac{K_{\text{total}}}{D_{\text{age}} + 1}
\end{equation}

\noindent with $K_{\text{total}}$ denoting cumulative karma and $D_{\text{age}}$ representing account age in days.

\subsection{Learning Algorithm Configuration}

Every platform's primary learner is an XGBoost gradient-boosted tree \cite{chen2016xgboost}. The initial parameter set appears in Table~\ref{tab:hyperparams}.

\begin{table}[htbp]
\caption{Starting XGBoost Parameter Set}
\label{tab:hyperparams}
\centering
\begin{tabular}{lcc}
\toprule
\textbf{Parameter} & \textbf{IG/RD/FB} & \textbf{TW/LI/YT} \\
\midrule
n\_estimators & 200 & 200 \\
max\_depth & 5 & 6 \\
learning\_rate & 0.1 & 0.1 \\
eval\_metric & logloss & logloss \\
random\_state & 42 & 42 \\
\bottomrule
\end{tabular}
\end{table}

To counteract label imbalance within the Facebook corpus, the minority-class weight is set as:

\begin{equation}
w_{\text{pos}} = \frac{|C_{\text{negative}}|}{|C_{\text{positive}}|}
\end{equation}

All classifiers undergo evaluation on a hold-out partition (80\% train, 20\% test, preserving original class proportions) and are additionally stress-tested via five rounds of cross-validation.

\subsection{Assembling the Voting Panel}

Accompanying each XGBoost learner are two algorithmically distinct partners:

\begin{itemize}
    \item \textbf{Bootstrap-Aggregated Forest:} 200 independent trees, maximum branching depth of 10, requiring at least 5 observations to justify a split.
    \item \textbf{Ridge-Penalized Linear Classifier:} L2 regularization with a 1{,}000-iteration convergence ceiling, preceded by z-score normalization of all inputs.
\end{itemize}

Per-fold accuracies for the supplementary algorithms are tabulated in Table~\ref{tab:ensemble_cv}.

\begin{table}[htbp]
\caption{Supplementary Model Accuracies Under Five-Fold Validation (\%)}
\label{tab:ensemble_cv}
\centering
\begin{tabular}{lccc}
\toprule
\textbf{Platform} & \textbf{XGB} & \textbf{Forest} & \textbf{Linear} \\
\midrule
Instagram & 98.8 & 98.3 $\pm$ 0.4 & 94.4 $\pm$ 1.2 \\
Twitter & 99.6 & 93.7 $\pm$ 11.6 & 94.0 $\pm$ 4.2 \\
Facebook & 96.0 & 85.0 $\pm$ 21.3 & 90.5 $\pm$ 13.5 \\
LinkedIn$^\dagger$ & 100.0 & 100.0 & 100.0 \\
YouTube$^\dagger$ & 99.9 & 100.0 & 100.0 \\
\bottomrule
\multicolumn{4}{l}{\footnotesize $^\dagger$Fabricated data; figures are non-generalizable.}
\end{tabular}
\end{table}

\subsection{Score Fusion Protocol}

The three per-model posteriors are merged via a weighted mean. Given feature vector $\mathbf{x}$, the composite score is:

\begin{equation}
\hat{P}(\mathbf{x}) = \frac{\sum_{k=1}^{3} w_k \cdot P_k(\mathbf{x})}{\sum_{k=1}^{3} w_k}
\end{equation}

\noindent where $P_k(\mathbf{x})$ denotes the $k$-th learner's estimated probability of inauthenticity and $w_k$ is proportional to that learner's ROC-AUC during validation. Consequently, models demonstrating stronger discriminative power wield correspondingly heavier influence on the final output.

Consensus among the trio is measured by the posterior spread:

\begin{equation}
\Delta = \max_k P_k(\mathbf{x}) - \min_k P_k(\mathbf{x})
\end{equation}

\noindent Values below 0.10 are tagged \textit{strong agreement}; values in $[0.10, 0.25)$ indicate \textit{partial agreement}; values at or above 0.25 signal \textit{disagreement}. When every model independently arrives at the same risk band (low/medium/high), a ``unanimous verdict'' badge is appended, offering the user an additional confidence cue.

\subsection{Automated Hyperparameter Refinement}

The four genuine-data XGBoost instances undergo Bayesian optimization via Optuna \cite{akiba2019optuna}, maximizing five-fold F1 across 50 sequential trials. The explored parameter landscape spans:

\begin{itemize}
    \item Tree count: $[100, 500]$; branching depth: $[3, 10]$
    \item Step magnitude: $[0.01, 0.3]$ sampled log-uniformly
    \item Leaf-weight floor: $[1, 10]$
    \item Row and column subsampling ratios: $[0.6, 1.0]$
    \item Complexity penalties: $\gamma$, $\alpha$ (L1), $\lambda$ (L2)
\end{itemize}

This procedure is deliberately skipped for the LinkedIn and YouTube classifiers: their non-overlapping synthetic distributions make any hyperparameter configuration equally effective, rendering the search meaningless.

\subsection{Attribution via Shapley Decomposition}

To make every verdict interpretable, we invoke the TreeExplainer kernel \cite{lundberg2017shap}, which recovers exact Shapley attributions for tree-based learners in polynomial time. The eight attributes carrying the largest absolute contributions are surfaced to the user. The decomposition satisfies:

\begin{equation}
f(\mathbf{x}) = \mathbb{E}[f(\mathbf{X})] + \sum_{i=1}^{M} \phi_i
\end{equation}

\noindent Here $\mathbb{E}[f(\mathbf{X})]$ is the baseline prediction averaged over the training population, $M$ is the feature count, and $\phi_i$ captures the marginal influence of the $i$-th attribute. A positive $\phi_i$ steers the output toward the ``inauthentic'' pole; a negative $\phi_i$ steers it toward ``genuine.''

% ============================================================================
% V. IMPLEMENTATION
% ============================================================================
\section{Implementation}

\subsection{Technology Choices}

Table~\ref{tab:techstack} itemizes the principal technologies employed.

\begin{table}[htbp]
\caption{Technology Portfolio}
\label{tab:techstack}
\centering
\begin{tabular}{ll}
\toprule
\textbf{Responsibility} & \textbf{Technology} \\
\midrule
Client Framework & React 19 + Vite 8 \\
Visual Language & Custom CSS (glassmorphic motifs) \\
Navigation & React Router DOM 7 \\
Data Visualization & Recharts 3.8 \\
Identity Provider & Firebase Auth (Email/Google/GitHub) \\
Document Store & Cloud Firestore \\
Server Framework & FastAPI + Uvicorn \\
Learning Algorithms & XGBoost + scikit-learn \\
Attribution Engine & SHAP TreeExplainer \\
Document Synthesis & fpdf2 \\
Async HTTP & httpx \\
Conversational AI & Groq Llama 3.3-70B + Google Gemini \\
In-Browser Scanner & Chrome Manifest V3 \\
Containerization & Docker + Docker Compose \\
Continuous Integration & GitHub Actions (4 stages) \\
Client-Side Testing & Vitest + Playwright \\
\bottomrule
\end{tabular}
\end{table}

\subsection{Endpoint Surface}

The server exposes over 25 routes partitioned into two authentication domains:

\textbf{Internal Surface (Firebase-protected):}
\begin{itemize}
    \item \textbf{Single Prediction} (\texttt{POST /api/predict/\{platform\}?ensemble=true}): Accepts platform-tailored attributes, returns fake-class probability, risk tier, SHAP attributions, and --- when ensemble mode is activated --- individual model scores with concordance analysis.
    \item \textbf{Bulk Prediction} (\texttt{POST /api/predict/batch/\{platform\}}): Ingests a list of records and returns per-record verdicts alongside aggregate statistics.
    \item \textbf{URL Inspector} (\texttt{POST /api/scan/url}): Identifies the target platform via URL pattern matching and, for Reddit, executes automated profile-data retrieval.
    \item \textbf{Document Generation}: Produces single-record PDF reports and multi-record ZIP archives, secured via tokenized download links with ten-minute expiry.
    \item \textbf{Conversational Proxy} (\texttt{POST /api/chat}): Routes user messages to Groq (primary) or Gemini (fallback), keeping all provider credentials server-side.
    \item \textbf{Evaluation Dashboard} (\texttt{GET /api/models/comparison}): Serves pre-computed accuracy, precision, recall, and AUC metrics for frontend visualization.
\end{itemize}

\textbf{Public Surface (API-key-protected):}
\begin{itemize}
    \item \textbf{Programmatic Analysis} (\texttt{POST /api/v1/analyze/\{platform\}}): Returns a trimmed JSON payload --- probability, risk tier, confidence --- without exposing SHAP internals or raw input echoes.
    \item \textbf{Credential Lifecycle}: Endpoints for key generation, masked listing, and revocation, with per-key request counters.
    \item \textbf{Platform Catalog} (\texttt{GET /api/v1/platforms}): Enumerates the six supported networks.
\end{itemize}

Input validation across every route is enforced by Pydantic v2 schemas with per-field boundary constraints, ensuring that out-of-range or malformed payloads are rejected with HTTP~422 before reaching any model.

\subsection{Defense-in-Depth Security Posture}

Multiple overlapping safeguards harden the API surface:

\begin{itemize}
    \item \textbf{Tiered middleware:} The outermost layer stamps each request with a unique X-Request-ID (honoring client-supplied IDs when present) and logs method, path, status, and latency. The middle layer enforces a sliding-window rate ceiling, persists its state to disk for crash resilience, and attaches X-RateLimit-Limit, X-RateLimit-Remaining, and X-RateLimit-Reset headers to every response. The innermost layer verifies Firebase ID tokens, gracefully degrading to an open bypass in development environments.
    \item \textbf{Strict input gates:} All 96+ numeric and categorical fields carry Pydantic \texttt{ge}/\texttt{le}/\texttt{min\_length}/\texttt{max\_length} annotations, blocking negative follower counts, out-of-range boolean flags, and empty batch lists at the schema boundary.
    \item \textbf{Environment isolation:} Secrets are loaded from \texttt{.env} files excluded via \texttt{.gitignore}; CORS origins are allow-listed (never wildcarded).
    \item \textbf{Dual-layer authentication:} Browser sessions present Firebase bearer tokens; programmatic clients present API keys.
    \item \textbf{Ephemeral caching:} Generated PDF reports live in memory under a ten-minute time-to-live eviction policy.
    \item \textbf{Test-safe mode:} Setting \texttt{TESTING=1} disables rate limiting during automated test runs without weakening production guards.
    \item \textbf{Client-side hardening:} The HTML entry point carries \texttt{X-Content-Type-Options}, \texttt{Referrer-Policy}, Open Graph, and Twitter Card meta tags.
\end{itemize}

\subsection{Visual Design}

The interface is rendered through over 3{,}700 lines of hand-authored CSS employing the following design language:

\begin{itemize}
    \item Theme duality via CSS custom properties, toggled with a 0.3-second crossfade, persisted to \texttt{localStorage}, and seeded from the operating system's color-scheme preference.
    \item Animated gradient backdrops with parallax-style motion.
    \item Frosted-glass card surfaces achieved through \texttt{backdrop-filter: blur} and semi-transparent borders.
    \item SVG-rendered probability gauges and gradient-filled ensemble-contribution bar charts.
    \item Platform-aware accent colors bound to CSS variables.
    \item Split-screen comparison layout with a feature-difference table and animated profile-swap transitions.
    \item Four-tab analytics dashboard: overview counters, temporal risk trends, model-performance radar plots, and threat heatmaps.
    \item Developer key-management pane with live usage counters and ready-to-copy cURL examples.
    \item Responsive grid scaffolding for cross-viewport adaptability.
    \item Floating chatbot widget with dual-provider failover (Groq $\to$ Gemini $\to$ offline fallback message).
\end{itemize}

\subsection{Quality Assurance}

Reliability is safeguarded through 110 automated checks distributed across three testing tiers:

\begin{table}[htbp]
\caption{Automated Test Inventory}
\label{tab:tests}
\centering
\begin{tabular}{lcc}
\toprule
\textbf{Module} & \textbf{Checks} & \textbf{Scope} \\
\midrule
\multicolumn{3}{l}{\textit{Server-side (pytest --- 80+ checks)}} \\
Prediction endpoints (6 + batch) & 20 & Correctness \\
SHAP structure \& data-source tags & 24 & Explainability \\
Schema boundary enforcement & 20 & Input safety \\
Public API key lifecycle & 12 & Auth \& routing \\
URL scanner \& evaluation & 9 & Detection logic \\
PDF/ZIP generation \& expiry & 8 & Reporting \\
Health diagnostics & 7 & Ops visibility \\
Chatbot proxy \& failover & 5 & Degradation \\
Rate-limit enforcement & 3 & Throttling \\
\midrule
\multicolumn{3}{l}{\textit{Client-side (Vitest --- 28 checks)}} \\
React component rendering & 28 & UI correctness \\
\midrule
\multicolumn{3}{l}{\textit{End-to-end (Playwright)}} \\
Full user journeys & --- & Integration \\
\midrule
\textbf{Aggregate} & \textbf{110+} & \textbf{All green} \\
\bottomrule
\end{tabular}
\end{table}

\subsection{Continuous Integration and Containerization}

A GitHub Actions pipeline orchestrates four sequential stages:

\begin{enumerate}
    \item \textbf{Server Tests:} Executes 80+ pytest checks under a \texttt{TESTING=1} guard.
    \item \textbf{Client Tests:} Runs Vitest unit checks followed by a production build to catch bundling regressions.
    \item \textbf{Container Build:} Verifies that both Docker Compose services assemble cleanly.
    \item \textbf{Integration Tests:} Boots both servers, installs Playwright chromium, and runs browser-level user-flow tests (triggered only on pushes to the main branch to conserve CI minutes).
\end{enumerate}

Docker Compose defines two services --- \texttt{backend} and \texttt{frontend} --- with health checks, restart policies, and a documented \texttt{--build-arg} override for injecting production API URLs at image-build time.

% ============================================================================
% VI. EXPERIMENTAL RESULTS
% ============================================================================
\section{Experimental Results}

\subsection{Hold-Out Classification Metrics}

Table~\ref{tab:results} reports performance on the stratified 80/20 test partitions.

\begin{table}[htbp]
\caption{Per-Platform Classification Metrics}
\label{tab:results}
\centering
\begin{tabular}{lccccc}
\toprule
\textbf{Platform} & \textbf{Acc.} & \textbf{Prec.} & \textbf{Rec.} & \textbf{F1} & \textbf{CV} \\
\midrule
Instagram & 98.6\% & 0.99 & 0.98 & 0.99 & 98.8\% \\
Twitter & 99.3\% & 1.00 & 0.99 & 0.99 & 99.6\% \\
Reddit & 94.0\% & 0.95 & 0.91 & 0.93 & 95.0\% \\
Facebook & 95.0\% & 0.83 & 0.90 & 0.86 & 96.0\% \\
LinkedIn$^\dagger$ & 100.0\% & 1.00 & 1.00 & 1.00 & 100.0\% \\
YouTube$^\dagger$ & 100.0\% & 1.00 & 1.00 & 1.00 & 99.9\% \\
\bottomrule
\multicolumn{6}{l}{\footnotesize $^\dagger$Synthetic training data; real-world accuracy will differ.} \\
\multicolumn{6}{l}{\footnotesize CV = five-fold cross-validation accuracy.}
\end{tabular}
\end{table}

The Instagram classifier delivers 98.6\% accuracy with tightly balanced precision (0.99) and recall (0.98). Twitter achieves the strongest genuine-data accuracy at 99.3\%, paired with an AUC of 0.9997. Reddit attains 94.0\% drawing on merely six features, underscoring the discriminative potency of karma-derived signals. Facebook reaches 95.0\% despite label imbalance, addressed through positive-class weighting.

\subsection{Algorithm Comparison and Ensemble Rationale}

To validate the ensemble design, all three learners were trained on identical splits. Table~\ref{tab:comparison} contrasts their standalone accuracies.

\begin{table}[htbp]
\caption{Standalone Accuracy: XGBoost vs.\ Random Forest vs.\ Logistic Regression}
\label{tab:comparison}
\centering
\begin{tabular}{lccc}
\toprule
\textbf{Platform} & \textbf{XGBoost} & \textbf{RF} & \textbf{LR} \\
\midrule
Instagram & 98.6\% & 98.3\% & 94.4\% \\
Twitter & 99.3\% & 93.7\% & 94.0\% \\
Facebook & 95.0\% & 85.0\% & 90.5\% \\
LinkedIn$^\dagger$ & 100.0\% & 100.0\% & 100.0\% \\
YouTube$^\dagger$ & 100.0\% & 100.0\% & 100.0\% \\
\bottomrule
\multicolumn{4}{l}{\footnotesize $^\dagger$Synthetic data; metrics non-indicative.}
\end{tabular}
\end{table}

XGBoost consistently leads on genuine-data benchmarks. The ensemble amalgamates all three through AUC-derived weights, delivering two practical advantages: (1) resilience against individual-model weaknesses and (2) inherent confidence signaling through concordance analysis. When all three classifiers concur on the risk tier, the system flags the outcome as ``unanimous,'' conveying heightened certainty. Divergent votes trigger an advisory prompt encouraging manual scrutiny.

Typical weight allocation assigns approximately 0.50 to XGBoost (strongest AUC), 0.30 to Random Forest (diversity via bagging), and 0.20 to Logistic Regression (diversity via linear boundaries).

\subsection{SHAP Feature Attributions}

Table~\ref{fig:shap} illustrates representative attributions for the Instagram classifier.

\begin{figure}[htbp]
\centerline{
\begin{tabular}{|l|r|}
\hline
\textbf{Attribute (Instagram)} & \textbf{$\phi$ Value} \\
\hline
Profile Picture Presence & $-2.8435$ \\
Follower / Followee Ratio & $-1.2140$ \\
Posts / Follower Density & $+0.8923$ \\
Completeness Index & $-0.7451$ \\
Biography Presence & $-0.5102$ \\
\hline
\end{tabular}
}
\caption{Top-5 SHAP attributions for the Instagram classifier. Negative values push the score toward ``genuine.''}
\label{fig:shap}
\end{figure}

Notable patterns across platforms:

\begin{itemize}
    \item \textbf{Instagram:} Avatar presence dominates, contributing a $+2.84$ shift toward ``fake'' when absent. The follower-to-followee ratio is the second-strongest signal.
    \item \textbf{Twitter:} Verification status and list-membership count emerge as the most separating dimensions, with verified badges producing strongly negative (genuine-leaning) attributions.
    \item \textbf{Reddit:} Daily karma velocity is the leading discriminator, effectively capturing the anomalous rapid-accumulation patterns characteristic of automated agents.
    \item \textbf{Facebook:} Community affiliation and account longevity dominate; spam accounts display markedly lower group participation.
\end{itemize}

\subsection{Sanity Checks on Constructed Profiles}

Table~\ref{tab:validation} shows model outputs for hand-crafted archetypes.

\begin{table}[htbp]
\caption{Prediction Outputs for Archetypal Profiles}
\label{tab:validation}
\centering
\begin{tabular}{llcc}
\toprule
\textbf{Platform} & \textbf{Archetype} & \textbf{$\hat{P}$} & \textbf{Tier} \\
\midrule
Instagram & Complete genuine & 0.00 & Low \\
Instagram & No avatar/bio & 1.00 & High \\
Twitter & Verified human & 0.01 & Low \\
Twitter & Minimal bot-like & 0.30 & Medium \\
Reddit & Veteran, high karma & 0.001 & Low \\
Reddit & New, negligible karma & 0.607 & High \\
Facebook & Engaged community member & 0.001 & Low \\
Facebook & URL-heavy poster & 0.216 & Low \\
\bottomrule
\end{tabular}
\end{table}

Instagram and Reddit exhibit strong class separation. Twitter differentiates clearly between verified humans and low-engagement bot patterns. Facebook's conservative bias is apparent: the URL-heavy archetype scores only 0.216, suggesting the model may benefit from additional spam-indicative features.

\subsection{Inference Speed}

Timing measurements on a consumer workstation (Intel i5, 16\,GB RAM):

\begin{itemize}
    \item Single prediction: $<$ 15\,ms
    \item SHAP computation: $<$ 50\,ms per instance
    \item Batch of 100 records: $<$ 2\,s
    \item PDF synthesis: $<$ 200\,ms
    \item End-to-end API round-trip (prediction + SHAP): $<$ 100\,ms
\end{itemize}

% ============================================================================
% VII. DISCUSSION
% ============================================================================
\section{Discussion}

\subsection{Advantages Over Prior Art}

\begin{enumerate}
    \item \textbf{Breadth of coverage.} Unlike single-network detectors, SocialGuard AI provides a consistent experience across six platforms under one architectural umbrella.
    \item \textbf{Ensemble-based confidence.} Fusing three algorithmically diverse learners through AUC-weighted voting not only improves robustness but also generates an organic confidence signal through concordance measurement --- a property absent from monolithic classifiers.
    \item \textbf{Transparent verdicts.} Per-prediction SHAP attributions let users and moderators trace the reasoning chain, fostering trust and enabling targeted investigation.
    \item \textbf{Deployment maturity.} Docker containerization, a four-stage CI pipeline, 110 automated checks across three testing tiers, dual authentication, persistent rate limiting with quota headers, and distributed request tracing collectively demonstrate production viability beyond a laboratory prototype.
    \item \textbf{Plug-in extensibility.} Supporting a new platform requires only defining a feature schema, training three classifiers, and registering an endpoint --- no changes to the ensemble engine, middleware, or frontend routing.
    \item \textbf{Open programmatic access.} The public API with key-gated authentication enables downstream integrations and automated workflows without exposing internal model details.
\end{enumerate}

\subsection{Acknowledged Limitations}

\begin{enumerate}
    \item \textbf{Synthetic training corpora (critical).} The LinkedIn and YouTube classifiers rest on artificially constructed data whose class distributions do not overlap. The perfect 100\% metrics are therefore artifacts of trivial separability, \textit{not} indicators of genuine discriminative ability. These classifiers exist solely to prove architectural extensibility and are flagged as ``Synthetic'' in every output channel --- the API (\texttt{data\_source: "synthetic"}), the interface, and the programmatic verification layer --- with automatic confidence attenuation applied.
    \item \textbf{Restricted platform APIs.} Only Reddit permits fully automated data retrieval; the remaining platforms require manual feature entry due to API access restrictions.
    \item \textbf{Modest corpus sizes.} Reddit (500 records) and Facebook (600 records) are small by modern standards, potentially limiting generalizability.
    \item \textbf{Concept drift.} Adversarial actors continuously refine evasion tactics; periodic retraining on fresh data is essential to maintain accuracy.
    \item \textbf{Partial ensemble coverage.} Reddit's continuous target labels precluded training supplementary Random Forest and Logistic Regression classifiers, limiting the ensemble to five of six platforms.
    \item \textbf{Volatile key storage.} Public API keys currently reside in memory and are lost on server restart; production deployments should migrate them to a durable store such as Firestore or Redis.
\end{enumerate}

\subsection{Prospective Extensions}

\begin{itemize}
    \item Sourcing authentic LinkedIn and YouTube corpora through institutional API partnerships or crowd-annotated data drives.
    \item Exploring cross-platform transfer learning to share latent representations between platform-specific classifiers.
    \item Incorporating natural-language content analysis --- comment toxicity scoring, posting-cadence anomaly detection --- as supplementary input dimensions.
    \item Deploying WebSocket-based streaming detection for continuous real-time monitoring.
    \item Replacing file-backed rate limiting and in-memory key storage with Redis to enable horizontal auto-scaling.
    \item Introducing adversarial robustness evaluations to quantify resilience against deliberate evasion strategies.
    \item Building model-versioning infrastructure with A/B traffic splitting for continuous model improvement.
\end{itemize}

% ============================================================================
% VIII. CONCLUSION
% ============================================================================
\section{Conclusion}

This paper introduced SocialGuard AI, an ensemble-driven detection framework spanning six social media platforms. By fusing XGBoost, Random Forest, and Logistic Regression through AUC-calibrated soft voting with concordance-based confidence measurement, the system achieves between 94.0\% and 99.3\% accuracy on genuine crowd-sourced datasets, with ROC-AUC values exceeding 0.97 across every real-data classifier.

The accompanying full-stack deployment --- a FastAPI server hardened with three-layer middleware (X-Request-ID tracing, X-RateLimit-* quota headers, Firebase token verification), a React 19 client with dual-theme rendering, a Chrome extension, dual-provider conversational assistance (Groq + Gemini), batch ingestion with PDF/ZIP reporting, and a key-gated public API --- validates that explainable, ensemble-powered inauthentic-account detection can operate as a production-grade web service. Docker Compose containerization and a four-stage GitHub Actions CI pipeline ensure repeatable builds. Dependability is attested by 110 automated checks spanning server-side unit tests, client-side component tests, and browser-level integration flows.

The pairing of SHAP attributions with ensemble concordance indicators distinguishes this work from opaque single-model classifiers, equipping users to understand both \textit{why} a profile is flagged and \textit{how confident} the system is in that judgment. While the synthetic LinkedIn and YouTube classifiers remain an acknowledged limitation, the modular architecture ensures that integrating genuine datasets --- when they become available --- requires no structural changes to the framework.

SocialGuard AI offers a comprehensive, transparent, and extensible platform that empowers individuals, researchers, and content moderators to confront the escalating challenge of inauthentic social media activity.

% ============================================================================
% REFERENCES
% ============================================================================
\begin{thebibliography}{00}

\bibitem{statista2026}
Statista, ``Number of social media users worldwide from 2017 to 2028,'' \textit{Statista Research Department}, 2026. [Online]. Available: https://www.statista.com/statistics/278414/number-of-worldwide-social-network-users/

\bibitem{cresci2020decade}
S. Cresci, ``A Decade of Social Bot Detection,'' \textit{Communications of the ACM}, vol. 63, no. 10, pp. 72--83, Oct. 2020.

\bibitem{varol2017}
O. Varol, E. Ferrara, C. A. Davis, F. Menczer, and A. Flammini, ``Online Human-Bot Interactions: Detection, Estimation, and Characterization,'' in \textit{Proc. AAAI ICWSM}, 2017, pp. 280--289.

\bibitem{ferrara2016rise}
E. Ferrara, O. Varol, C. Davis, F. Menczer, and A. Flammini, ``The Rise of Social Bots,'' \textit{Communications of the ACM}, vol. 59, no. 7, pp. 96--104, Jul. 2016.

\bibitem{yang2020scalable}
K.-C. Yang, O. Varol, P.-M. Hui, and F. Menczer, ``Scalable and Generalizable Social Bot Detection through Data Selection,'' in \textit{Proc. AAAI}, 2020, pp. 1096--1103.

\bibitem{akyon2019instagram}
F. C. Akyon and E. Kalfaoglu, ``Instagram Fake and Automated Account Detection,'' in \textit{Proc. IEEE INISTA}, 2019, pp. 1--7.

\bibitem{fire2014friend}
M. Fire, D. Kagan, A. Elyashar, and Y. Elovici, ``Friend or Foe? Fake Profile Identification in Online Social Networks,'' \textit{Social Network Analysis and Mining}, vol. 4, no. 1, pp. 1--23, 2014.

\bibitem{boshmaf2011socialbot}
Y. Boshmaf, I. Muslukhov, K. Beznosov, and M. Ripeanu, ``The Socialbot Network: When Bots Socialize for Fame and Money,'' in \textit{Proc. ACSAC}, 2011, pp. 93--102.

\bibitem{hurtado2019bot}
S. Hurtado, P. Ray, and R. Marculescu, ``Bot Detection in Reddit Political Discussion,'' in \textit{Proc. SNS}, 2019, pp. 30--35.

\bibitem{kudugunta2018deep}
S. Kudugunta and E. Ferrara, ``Deep Neural Networks for Bot Detection,'' \textit{Information Sciences}, vol. 467, pp. 312--322, Oct. 2018.

\bibitem{lundberg2017shap}
S. M. Lundberg and S.-I. Lee, ``A Unified Approach to Interpreting Model Predictions,'' in \textit{Advances in Neural Information Processing Systems (NeurIPS)}, 2017, pp. 4765--4774.

\bibitem{rathore2023xai}
H. Rathore, C. Samavedhi, and S. K. Gupta, ``XAI-based Network Intrusion Detection using SHAP and LIME,'' in \textit{Proc. IEEE ISDFS}, 2023, pp. 1--6.

\bibitem{mohawesh2023explainability}
R. Mohawesh et al., ``Explainability in AI-Based Content Moderation: A Comprehensive Survey,'' \textit{Artificial Intelligence Review}, vol. 56, pp. 10145--10192, 2023.

\bibitem{chen2016xgboost}
T. Chen and C. Guestrin, ``XGBoost: A Scalable Tree Boosting System,'' in \textit{Proc. ACM KDD}, 2016, pp. 785--794.

\bibitem{zhang2019fraud}
Z. Zhang, X. Zhou, X. Zhang, L. Wang, and P. Wang, ``A Model Based on Convolutional Neural Network for Online Transaction Fraud Detection,'' \textit{Security and Communication Networks}, vol. 2018, pp. 1--9, 2019.

\bibitem{dhaliwal2018effective}
S. S. Dhaliwal, A.-A. Nahid, and R. Abbas, ``Effective Intrusion Detection System Using XGBoost,'' \textit{Information}, vol. 9, no. 7, p. 149, 2018.

\bibitem{akiba2019optuna}
T. Akiba, S. Sano, T. Yanase, T. Ohta, and M. Koyama, ``Optuna: A Next-generation Hyperparameter Optimization Framework,'' in \textit{Proc. ACM KDD}, 2019, pp. 2623--2631.

\bibitem{instagram_dataset}
``Instagram Fake Profile / Automated Account Detection,'' Kaggle, 2022. [Online]. Available: https://www.kaggle.com/datasets/rajrohan/instagram-fake-spammer-genuine-accounts

\bibitem{twitter_dataset}
``Twitter Bot Detection Dataset (Genuine and Fake User Accounts),'' Kaggle, 2021. [Online]. Available: https://www.kaggle.com/datasets/davidmartngutirrez/twitter-bots-accounts

\bibitem{reddit_dataset}
``Reddit Dead Internet Analysis Dataset,'' Custom Collection, 2026.

\bibitem{facebook_dataset}
``Facebook Spam Detection Dataset,'' Kaggle, 2022. [Online]. Available: https://www.kaggle.com/datasets/mfaisalqureshi/facebook-spam-dataset

\end{thebibliography}

\end{document}
